\subsection{INTEGRAL VARIANT OF THE POINCAR\'E-BIRKHOFF-WITT THEOREM}

Let g be a finite dimensional Q-Lie algebra, \(U(\mathfrak{g})\) its enveloping bigebra. If I is a totally ordered set, \(\mathbf{x} = (x_i)_{i \in \mathrm{I}}\) a family of elements of g, and \(\mathbf{n} = (n_i)_{i \in \mathrm{I}} \in \mathbf{N}^{(\mathrm{I})}\) a multi-index, put

\[
\mathbf {x} ^ {(\mathbf {n})} = \prod_ {i \in \mathrm{I}} \frac {x _ {i} ^ {n _ {i}}}{n _ {i} !},\tag{5}
\]

the product being calculated in \(\mathrm{U}(\mathfrak{g})\) in accordance with the ordered set I.

THEOREM 1. Let \(\mathcal{U}\) be a biorder in the bigebra \(\mathrm{U}(\mathfrak{g})\) . Let \(\mathcal{G} = \mathcal{U} \cap \mathfrak{g}\) , which is an order in the Lie algebra \(\mathfrak{g}\) . Let \((x_i)_{i \in \mathrm{I}}\) be a basis of \(\mathcal{G}\) . Give I a total order, and assume that we are given, for all \(\mathbf{n} \in \mathbf{N}^{\mathrm{I}}\) , an element {[}n{]} of \(\mathcal{U}\) such that {[}n{]} - \(\mathbf{x}^{(\mathbf{n})}\) has filtration \(< |\mathbf{n}|\) in \(\mathrm{U}(\mathfrak{g})\) . Then, the family of the {[}n{]} for \(\mathbf{n} \in \mathbf{N}^{\mathrm{I}}\) is a basis of the Z-module \(\mathcal{U}\) .

For \(p \in N\) , let \(U_{p}(\mathfrak{g})\) be the set of elements of \(U(\mathfrak{g})\) of filtration \(\leq p\) ; then the images in \(U_{p}(\mathfrak{g}) / U_{p-1}(\mathfrak{g})\) of the \(\mathbf{x}^{(\mathbf{n})}\) such that \(|n| = p\) form a basis of this Q-vector space (Chap. I, §2, no. 7, Th. 1); hence the {[}n{]} form a basis of the Q-vector space \(U(\mathfrak{g})\) . It remains to prove the following assertion (in which we put \(M = N^{I}\) ):

\((^{*})\) if \(u\in \mathcal{U}\) , \((a_{\mathbf{n}})\in \mathbf{Z}^{(\mathrm{M})}\) , and \(d\in \mathbf{N} - \{0\}\) are such that

\[
d u = \sum_ {\mathbf {n} \in \mathrm{M}} a _ {\mathbf {n}} [ \mathbf {n} ],\tag{6}
\]

then \(d\) divides each \(a_{\mathbf{n}}\) .

For each integer \(r \geq 0\) , introduce the iterated coproduct

\[
c _ {i}: \mathcal {U} \rightarrow \mathbf {T} ^ {r} (\mathcal {U}) = \mathcal {U} \otimes \mathcal {U} \otimes \dots \otimes \mathcal {U};
\]

by definition, \(c_0\) is the counit of \(\mathcal{U}\) , \(c_1 = \mathrm{Id}_{\mathcal{U}}\) , \(c_2 = c\) (the coproduct of \(\mathcal{U}\) ), and, for \(r \geq 2\) , \(c_{r+1}\) is defined as the composite \(p \circ (c_r \otimes 1) \circ c\) :

\[
\mathcal {U} \stackrel {{c}} {{\longrightarrow}} \mathcal {U} \otimes_ {\mathbf {Z}} \mathcal {U} \stackrel {{c _ {r} \otimes 1}} {{\longrightarrow}} \mathbf {T} ^ {r} (\mathcal {U}) \otimes_ {\mathbf {Z}} \mathcal {U} \stackrel {{p}} {{\longrightarrow}} \mathbf {T} ^ {r + 1} (\mathcal {U})
\]

where p is defined by using the multiplication in the algebra \(\mathbf{T}(\mathcal{U})\) . Further, consider the canonical projection \(\pi\) of U onto \(U^{+} = Ker c_{0}\) , and the composite

\[
c _ {r} ^ {+} = \mathbf {T} ^ {r} (\pi) \circ c _ {r}: \mathcal {U} \to \mathbf {T} ^ {r} (\mathcal {U} ^ {+}).
\]

Lemma 1. Let \(\mathbf{n} \in \mathbf{N}^{\mathrm{I}}\) . If \(|\mathbf{n}| < r\) , then \(c_{r}^{+}([ \mathbf{n}]) = 0\) . If \(|\mathbf{n}| = r\) , then

\[
c _ {r} ^ {+} ([ \mathbf {n} ]) = \sum_ {\varphi} x _ {\varphi (1)} \otimes x _ {\varphi (2)} \otimes \dots \otimes x _ {\varphi (r)},\tag{7}
\]

where \(\varphi\) belongs to the set of maps from \(\{1,2,\ldots,r\}\) to I which take each value \(i\in I\) \(n_{i}\) times.

By Prop. 1,

\[
c _ {r} \left(\mathbf {x} ^ {(\mathbf {n})}\right) = \sum \mathbf {x} ^ {\left(\mathbf {p} _ {1}\right)} \otimes \dots \otimes \mathbf {x} ^ {\left(\mathbf {p} _ {r}\right)}
\]

where the summation extends over the set of sequences \((\mathbf{p}_{1},\ldots,\mathbf{p}_{r})\) of r elements of M such that \(p_{1}+\cdots+p_{r}=n\) . In view of Chap. II, §1, no. 3, Prop. 6, the map \(c_{r}^{+}\) , extended by linearity to a map from \(\mathrm{U}(\mathfrak{g})\) to \(\mathbf{T}^{r}(\mathrm{U}^{+}(\mathfrak{g}))\) , vanishes on \(\mathrm{U}_{r-1}(\mathfrak{g})\) . It follows that, for \(r\geq|\mathbf{n}|\) ,

\[
c _ {r} ^ {+} ([ \mathbf {n} ]) = c _ {r} ^ {+} (\mathbf {x} ^ {(\mathbf {n})}) = \sum \pi (\mathbf {x} ^ {(\mathbf {p} _ {1})}) \otimes \dots \otimes \pi (\mathbf {x} ^ {(\mathbf {p} _ {r})}).\tag{8}
\]

For \(r > |n|\) , the relation \(p_{1} + \cdots + p_{r} = n\) implies that at least one of the \(p_{i}\) is zero, so \(c_{r}^{+}([n]) = 0\) . For \(r = |n|\) , the only non-zero terms of the third member of (8) are those for which \(|p_{1}| = \cdots = |p_{r}| = 1\) , hence (7).

We return to the proof of Th. 1. We retain the notations of (*) and prove, by descending induction on \(|n|\) , that d divides \(a_{n}\) , which is clear when \(|n|\) is sufficiently large. If d divides \(a_{n}\) for \(|n| > r\) then, putting

\[
u ^ {\prime} = u - \sum_ {| \mathbf {n} | > r} (a _ {\mathbf {n}} / d) [ \mathbf {n} ] \in \mathcal {U},
\]

we have

\[
d u ^ {\prime} = \sum_ {| \mathbf {n} | \leq r} a _ {\mathbf {n}} [ \mathbf {n} ].\tag{9}
\]

For any map \(\varphi\) from \(\{1,2,\ldots,r\}\) to I, put

\[
e _ {\varphi} = x _ {\varphi (1)} \otimes \dots \otimes x _ {\varphi (r)}
\]

and \(a_{\varphi} = a_{\mathbf{n}}\) where \(\mathbf{n} = (\operatorname{Card}\varphi^{-1}(i))_{i\in \mathrm{I}}\) . By Lemma 1, (9) implies that

\[
d c _ {r} ^ {+} (u ^ {\prime}) = \sum_ {\varphi \in \mathrm{I} ^ {r}} a _ {\varphi} e _ {\varphi}\tag{10}
\]

so \(c_r^+(u') \in \mathbf{T}^r(\mathcal{U}^+) \cap \mathbf{Q}\mathbf{T}^r(\mathcal{G})\) . But the submodule \(\mathcal{G}\) of \(\mathcal{U}^+\) is a direct factor (Algebra, Chap. VII, §4, no. 3, Cor. of Th. 1), so the submodule \(\mathbf{T}^r(\mathcal{G})\) is a direct factor of \(\mathbf{T}^r(\mathcal{U}^+)\) , and hence \(c_r^+(u') \in \mathbf{T}^r(\mathcal{G})\) . On the other hand, the \(x_i\) form a basis of \(\mathcal{G}\) by hypothesis, so the \(e_\varphi\) form a basis of \(\mathbf{T}^r(\mathcal{G})\) . Then (10) proves that \(d\) divides the \(a_\varphi\) , that is, the \(a_{\mathbf{n}}\) for \(|\mathbf{n}| = r\) . This proves (*).
% semantic-fragments: legacy-input-seam
\relax{}
